El objetivo de este artículo es analizar la integración de frameworks backend y patrones de diseño, específicamente el Modelo-Vista-Controlador (MVC), en el desarrollo de servicios web. Se destaca la importancia de los servicios RESTful en la comunicación entre cliente y servidor, como una alternativa eficiente a tecnologías legadas, utilizando métodos HTTP (GET, POST, PUT, DELETE). El estudio se centra en la comparación de dos frameworks principales: Spring Boot (Java) y Node.js (JavaScript), reconociendo sus fortalezas distintivas: Spring Boot por su robustez, seguridad y orientación empresarial, y Node.js por su rapidez, asincronía y escalabilidad en aplicaciones que requieren baja latencia. La integración del patrón MVC es crucial para la organización, mantenibilidad y testabilidad del código, ya que separa la lógica de negocio (Modelo), la interfaz (Vista) y el flujo de control (Controlador). Se concluye que la elección del framework backend (Spring Boot o Node.js) debe estar alineada con el contexto del proyecto, las necesidades de seguridad y la madurez del equipo, demostrando que la combinación de arquitecturas sólidas y herramientas adecuadas, gestionadas bajo la Metodología Scrum, es fundamental para el éxito del software.
